\section{Simple types: substitution, reduction, and termination}

\begin{annexitem}{annex:p7.stlc.substitution}{thm:p07-stlc-substitution}{Simultaneous typing substitution}
\begin{proof}
First establish the context facts used in the induction. Enlarging a context
by fresh declarations preserves a variable judgment by membership. It
preserves an application by applying the induction hypothesis to its two
premises. At an abstraction, rename the bound variable away from the added
declarations, apply the induction hypothesis to its body, and reapply the
abstraction rule. Reordering a context has the same three cases: membership
is unchanged, application uses its two premises, and abstraction retains
the fresh final declaration for its body. The renaming used here is justified
by the identical induction: replace the designated variable in the membership
case, propagate through an application, and rename an intervening binder
fresh before propagating through an abstraction. Thus the argument does not
assume a context rule absent from the registered calculus.

Now induct on the given derivation of \(\Theta\vdash M:\alpha\), using
the abbreviated turnstile only in this proof. If its last rule is a variable
rule, \(M=x\), the required derivation is exactly the hypothesis for
\(\rho(x)\). If it is application, the premises derive
\(P:\beta\mathbin{\to_{\mathrm{ty}}}\alpha\) and \(Q:\beta\).
Induction gives both substituted judgments in \(\Delta\), and application
derives \(P\rho\,Q\rho=(PQ)\rho:\alpha\).

For an abstraction, write \(M=\lambda y{:}\beta.P\), with premise
\(\Theta,y:\beta\vdash P:\gamma\) and
\(\alpha=\beta\mathbin{\to_{\mathrm{ty}}}\gamma\).
Choose \(y\) fresh for \(\Delta\), the substitution domain, and the free
variables of every image. In the target context \(\Delta,y:\beta\), the
old images retain their types by weakening, and the additional substitution
image of \(y\) is \(y\) itself. Apply induction with this extended
substitution. The result is \(\Delta,y:\beta\vdash P\rho:\gamma\).
Abstraction discharges the added declaration and gives
\(\Delta\vdash\lambda y{:}\beta.P\rho:\alpha\), exactly the
capture-avoiding substitution into \(M\). These exhaust the typing rules.
For the single-variable consequence, use the identity images for the
variables of \(\Theta\), and the specified image \(N\) for \(x\).
\end{proof}
\end{annexitem}

\begin{annexitem}{annex:p7.stlc.subject-reduction}{thm:p07-stlc-subject-reduction}{Subject reduction for simple types}
\begin{proof}
For a contraction at the root, typing inversion for application and
abstraction gives
\[
 \Theta,x:\beta\vdash_{\mathsf{STLC}}P:\alpha,
 \qquad \Theta\vdash_{\mathsf{STLC}}Q:\beta
\]
for \((\lambda x{:}\beta.P)Q\). The substitution lemma gives the
judgment \(\Theta\vdash_{\mathsf{STLC}}P[Q/x]:\alpha\) for its
contractum. Inversion is valid because the last typing rule is determined
by the outer constructor, and equality of arrow types requires equality of
their argument and result types.

For an arbitrary position, induct on the number of enclosing constructors.
There is no contraction inside a variable. Inside an abstraction, apply
induction in its context extended by the bound variable and reapply the
abstraction rule. Inside the function of an application, induction preserves
the function's entire arrow type while the argument judgment is unchanged.
Inside its argument, induction preserves the argument type while the function
judgment is unchanged. Reapply application in both cases. Root, abstraction,
and these two application cases are every clause of compatible reduction.
Iteration proves the finite-reduction assertion, including the empty sequence.
\end{proof}
\end{annexitem}

\begin{annexitem}{annex:p7.stlc.strong-normalization}{thm:p07-stlc-strong-normalization}{Strong normalization by reducibility}
\begin{proof}
Use the raw-term predicates of Definition~\ref{def:p07-stlc-reducibility}.
All term equalities in this proof respect alpha-equivalence. There are
finitely many redex positions in a finite term, and contracting a position
determines one reduct modulo alpha-equivalence. For a strongly normalizing
term \(M\), let \(h(M)\) be the greatest length of a reduction starting
at \(M\). This is a finite natural number. Indeed, the reduction tree is
finitely branching. If it had paths of unbounded length, some immediate
successor would also have unbounded remaining path lengths; repeatedly choose
the first such successor in the finite position order. This constructs an
infinite reduction, a contradiction. Consequently
\[
 M\to_\beta M'\quad\Longrightarrow\quad h(M')<h(M)
 \qquad(M\in\mathsf{SN}_\beta).
\]
Every reduct and every subterm of a strongly normalizing term is strongly
normalizing: an infinite sequence from either would extend to, or lift into,
one from the original term.

\paragraph{Substitution and reduction.}
We will use two elementary facts on raw terms:
\[
 P\to_\beta P'\ \Longrightarrow\ P\rho\to_\beta P'\rho,
 \qquad
 U\to_\beta U'\ \Longrightarrow\ P[U/x]\to_\beta^*P[U'/x].
 \tag{S}
\]
For the first, rename the binder of a contracted redex fresh for the
substitution and its images. Substituting into
\((\lambda y{:}\gamma.A)B\) leaves a beta redex with contractum
\(A\rho[B\rho/y]\). This is alpha-equal to \((A[B/y])\rho\).
The equality follows by induction on \(A\): variables distinguish \(y\)
from the other variables, applications distribute both substitutions to
both children, and abstractions use a binder fresh for both substitutions.
Induction on the context surrounding the redex proves the first fact in
all positions, with a genuine contraction still present after substitution.
For the second fact, reduce the designated occurrence inside each copy of
\(U\) inserted for a free occurrence of \(x\). There are finitely many
such disjoint copies. If there are none, use zero steps. Formally, induction
on \(P\) uses the given reduction at \(P=x\), reflexivity at other
variables, concatenation inside the two application positions, and reduction
inside a freshly renamed abstraction. The same facts apply to finite
sequences by iteration.

\paragraph{The three closure properties.}
For every type \(\alpha\), prove simultaneously:
\begin{enumerate}
 \item if \(M\in\mathcal R_\alpha\), then \(M\in\mathsf{SN}_\beta\);
 \item if \(M\in\mathcal R_\alpha\) and \(M\to_\beta M'\), then
       \(M'\in\mathcal R_\alpha\);
 \item if \(M\) is neutral and every immediate reduct of \(M\) is in
       \(\mathcal R_\alpha\), then \(M\in\mathcal R_\alpha\).
\end{enumerate}
At a base type the first is the definition. The second holds because an
infinite sequence from a successor would extend to one from \(M\).
For the third, any infinite sequence must begin with an immediate successor,
contradicting that successor's strong normalization. In particular, a variable,
being neutral with no reducts, belongs to the base predicate.

Let the type be \(\alpha\mathbin{\to_{\mathrm{ty}}}\beta\), and suppose
all three properties proved at its proper subtypes. To prove the first, take
a variable \(z\). Property 3 at \(\alpha\) puts \(z\) in
\(\mathcal R_\alpha\), so \(Mz\in\mathcal R_\beta\). Property 1 at
\(\beta\) makes \(Mz\) strongly normalizing. An infinite reduction of
\(M\) would remain an infinite reduction in the function position of
\(Mz\), so \(M\) is strongly normalizing.

For property 2, let \(U\in\mathcal R_\alpha\). Then
\(MU\in\mathcal R_\beta\) and \(MU\to_\beta M'U\).
Property 2 at \(\beta\) puts \(M'U\) in \(\mathcal R_\beta\).
As this holds for every such \(U\), the arrow definition gives the result.

For property 3, fix neutral \(M\) whose immediate reducts lie in the arrow
predicate. For every \(U\in\mathcal R_\alpha\), prove
\(MU\in\mathcal R_\beta\) by induction on \(h(U)\), which exists by
property 1 at \(\alpha\). A reduct of \(MU\) is either \(M'U\), which
is in \(\mathcal R_\beta\) by the assumption on \(M'\), or \(MU'\).
In the latter case property 2 at \(\alpha\) gives
\(U'\in\mathcal R_\alpha\), and the smaller height allows the induction
hypothesis. There is no root beta contraction, since \(M\) is not an
abstraction. The application \(MU\) is neutral, so property 3 at \(\beta\)
places it in \(\mathcal R_\beta\). This proves property 3 at the arrow
type and completes the induction on types. Variables therefore belong to
\(\mathcal R_\alpha\) at every type. Property 2 also holds for finite
reductions, by iteration.

\paragraph{Controlled beta expansion.}
For any type \(\beta\), if \(P,U\in\mathsf{SN}_\beta\) and
\(P[U/x]\in\mathcal R_\beta\), then
\[
       (\lambda x{:}\alpha.P)U\in\mathcal R_\beta.       \tag{B}
\]
Here the annotation \(\alpha\) is any fixed simple type; raw-term membership
does not presume a typing derivation. Induct on \(h(P)+h(U)\).
The root reduct is \(P[U/x]\), in the required predicate by assumption.
A reduction in \(P\) gives \((\lambda x{:}\alpha.P')U\).
The first fact (S) and property 2 give
\(P'[U/x]\in\mathcal R_\beta\); the strictly smaller height permits
induction. A reduction in \(U\) gives \((\lambda x{:}\alpha.P)U'\).
The second fact (S), including its zero-copy case, and finite reduction
closure give \(P[U'/x]\in\mathcal R_\beta\); induction again applies.
These are all immediate reducts. The entire application is neutral, so
property 3 places it in \(\mathcal R_\beta\), proving (B).

It follows that if \(P[U/x]\in\mathcal R_\beta\) for every
\(U\in\mathcal R_\alpha\), then
\(\lambda x{:}\alpha.P\in
 \mathcal R_{\alpha\mathbin{\to_{\mathrm{ty}}}\beta}\).
To justify the strong-normalization premise needed in (B), choose a variable
\(z\) fresh for \(P\). It belongs to \(\mathcal R_\alpha\), so
\(P[z/x]\in\mathcal R_\beta\subseteq\mathsf{SN}_\beta\).
By the first fact (S), any infinite reduction of \(P\) would give an
infinite reduction of \(P[z/x]\). Thus \(P\) is strongly normalizing.
Every \(U\in\mathcal R_\alpha\) is strongly normalizing by property 1,
and (B) now applies separately to every such argument. This proves the
abstraction property without assuming that an arbitrary beta expansion
preserves strong normalization.

\paragraph{Fundamental lemma.}
If \(\Theta\vdash_{\mathsf{STLC}}M:\alpha\) and a simultaneous raw
substitution \(\rho\) sends every \(x:\gamma\) in \(\Theta\) to a
member of \(\mathcal R_\gamma\), then \(M\rho\in\mathcal R_\alpha\).
Induct on the typing derivation. At a variable the assertion is its
substitution hypothesis. At application, the two induction hypotheses
give membership in the arrow and argument predicates; their application
belongs to the result predicate by its definition.

At an abstraction \(\lambda y{:}\gamma.P\), first rename \(y\) fresh
for the substitution domain and all images. For any
\(U\in\mathcal R_\gamma\), extend the substitution with \(y\mapsto U\).
The induction hypothesis for the premise gives
\(P\rho[U/y]\in\mathcal R_\delta\), where \(\delta\) is the body type.
The equality with the extended simultaneous substitution follows from the
freshness condition; induction on the body gives the variable, application,
and freshly renamed abstraction cases. The abstraction property just proved
therefore places \(\lambda y{:}\gamma.P\rho\) in
\(\mathcal R_{\gamma\mathbin{\to_{\mathrm{ty}}}\delta}\).
This is the required substituted abstraction. All typing cases are covered.

Finally choose the identity substitution on the finite context. Every
variable belongs to the predicate indexed by its assigned type, so the
fundamental lemma gives \(M\in\mathcal R_\alpha\). Property 1 yields
strong normalization, for open as well as closed terms. A maximal reduction
cannot be infinite and cannot stop at a reducible term by its definition.
It therefore ends in a normal term, whose type is preserved by
Theorem~\ref{thm:p07-stlc-subject-reduction}.
\end{proof}
\end{annexitem}

\begin{annexitem}{annex:p7.stlc.inert-constants}{thm:p07-stlc-inert-constants}{Typed inert constants}
\begin{proof}
A finite input term uses only finitely many constants. Replace each different
constant occurring in it by a different fresh variable assigned that constant's
declared simple type. Choose these variables fresh for every variable in the
term and its typing context. Induction on the extended typing derivation gives
a derivation of the replaced term in the ordinary calculus with these fresh
declarations added: a constant leaf becomes a variable leaf; original variable
leaves use weakening; application and abstraction retain their rules, with
binders renamed fresh when necessary.

The replacement commutes with capture-avoiding substitution, by induction on
the substituted term. Hence every beta contraction in the original term gives
one beta contraction at the same position in the replaced term. Constants
are inert, so there is no further reduction case to consider. An infinite
beta sequence of the original term would thus produce an infinite sequence
from a term typable in the ordinary calculus, contradicting strong
normalization. The reasoning uses only the finitely many constants in the
input. Beta reduction can duplicate or erase them but cannot introduce a
new constant, so an infinite constant signature causes no additional case.
\end{proof}
\end{annexitem}
